% 原书第 371 页（PDF 第 384 页）
\BookPage{371}{384}
\BookHeading{年年红药为谁开—温史话扬州}{年年红药为谁开\\[3pt]—温史话扬州}{冯尔健*}
「维桑与梓，必恭敬止」。身为邑人，自然对故乡扬州心怀崇敬并衷心希望它从来都是城泰民安、社会和谐，居者怡而远者慕，像古人诗句「人生只爱扬州住」「人生只合扬州死」所咏赞的那样美好可人。然而，古扬州的实际情况究竟如何？这种桑梓情怀与赞美诗句会不会只是一厢情愿、理想化的浪漫？

要了解古扬州实际上是怎样的城市，当然主要只能凭借来自古人笔下的相关文史资料。遗憾的是这类资料不仅数量有限，而且其所反映的城市面貌与民众生活又受限于操笔者本人的主客观因素，难以提供比较完整的全息图像。

资料不足，尤须解读得当。古人笔下的扬州，无论见仁见智、见善见恶、见美见丑，作为前车之鉴，于今人建设和谐社会皆有一定参考价值。故治史者应力求全其原貌展示于世，不宜断章取义、讳恶蔽丑，尤忌饰丑为美，误导后生。

笔者对活跃于清朝乾隆时期若干乡贤（汪中、郑燮、李斗等）有关扬州的描述粗加梳理之后，脑海中清晰地展现出两幅似即似离的图景：一幅是英雄辈出、不负天下、利尽四海、繁荣昌盛、百业兴隆、帝王巡幸、诗酒文会、十里长街、嘉园丽舍、朱楼翠幔、笙歌彻晓、风雅时尚…猗欤休哉！犹如一座座光艳夺目、金碧辉煌的宝塔！而在它们的四周背景上呈现的则是另一幅图景:灰暗，模糊，茫茫中隐隐可见萧娘的泪眼、桃叶的愁眉、栖身寒寺的贫士、鸣咽夜哭的荒冢…。两幅图景互相混杂，无论持何种视角欣赏，都难称谐美。本文试图借助历代有关扬州的史籍与

\vfill
\begin{BookNotes}
\BookNote{*}{作者任职于中国科学院电工研究所。}
\BookNote{1.}{《诗经·小弁》。}
\BookNote{2.}{（清）黄慎，《维扬竹枝词》。}
\BookNote{3.}{（唐）张祜，〈纵游淮南)。}
\end{BookNotes}

% 原书第 372 页（PDF 第 385 页）
\BookPage{372}{385}
\noindent 诗文〈芜城赋〉4、〈扬州慢〉5、〈广陵对〉6等对其不谐之处稍加辨析，以期于扬州今后之谐美有所裨益。

\BookSubheading{（一）英雄之城的背后}

我们先从乾隆时期乡贤汪中笔下的古扬州说起。

汪中（1745-1794）于乾隆五十二年（1787），应浙江学政朱珪之询，面述历朝历代「广陵之事」，并据以撰文〈广陵对〉，对拥有二千余年历史的乡邦广陵城作如下之生动描述：

\begin{BookQuote}
以广陵一城之地，天下无事，则煮海为盐，使万民食其业，上输少府，以宽农亩之力，及川渠所转，百货通焉，利尽四海；一旦（天下）有变，进则翼载天子，立桓文之功，退则保据州土、力谋恢复，不幸天长丧乱，知勇俱困，犹复与民守之，效死勿去，以明为人臣之义，历十有八姓、二千余年，而亡城降子，不出于其间。
\end{BookQuote}

总评充满自豪、掷地有声：

\begin{BookQuote}
由是言之，广陵何负于天下哉！
\end{BookQuote}

显见，汪中笔下所展示的广陵是一座英雄城。朱珪不仅产生强烈共鸣，而且评价比他更高：

\begin{BookQuote}
卓哉言乎！…天下百郡，洵无若广陵者。后之过者，式其城焉，可也！
\end{BookQuote}

朱珪是清朝嘉庆帝的师傅，熟谂史籍，判定就「忠君死节」事迹而言，天下百郡确实无一能及扬州，故而希望过往客向它致敬！

<广陵对〉所胪列的历代「死节」人士中，包括抗元之李庭芝与抗清之史可法，他们为抵御外族侵略率领军民坚守孤城扬州，最后为国壮烈牺牲，足为扬州增

\vfill
\begin{BookNotes}
\BookNote{4.}{《昭明文选》，卷2。}
\BookNote{5.}{《宋词选》，上海古籍，1978。}
\BookNote{6.}{《汪容甫文笺》，人民文学出版社，1958。}
\end{BookNotes}

% 原书第 373 页（PDF 第 386 页）
\BookPage{373}{386}
\noindent 色，值得邑人引以自豪。

尤其可贵的是，李庭芝在南宋全国土崩瓦解、皇太后与皇帝被俘北上并一再降旨令降的情况下，仍然坚持抗元斗争直至最后捐躯：7

\begin{BookQuote}
宋亡，谢太后及瀛国公（赵显）为诏，谕之降。庭芝登城曰：「奉诏守城，未闻有诏谕降也！」已而，两宫入朝（被掳北上），至瓜州，复诏庭芝曰：「比诏卿纳款，日久未报，岂未悉吾意，尚欲固圉邪？今吾与嗣君既已臣服，卿尚为谁守之？」庭芝不答，令发弩射使者，毙一人，余皆退去。姜才出兵夺两宫，不克，复闭城守。…（庭芝）被执，…死之日，扬之民皆涕下。
\end{BookQuote}

显见，李庭芝之死与其说是为一姓而死节，毋宁说是为民族而死难，其价值远在汪、朱传统忠君观念（无须赘言）之上，值得我们中华民族永远缅怀与崇敬。

<广陵对〉通篇列举的英雄人物多达二十以上，一一随人述事，轰轰烈烈，约占全文80\%以上篇幅。惟其中无一言语及「与民守之，效死勿去」之民，不能不说是汪文的一大缺笔。且看文中有关李庭芝与史可法是怎样写的：

\begin{BookQuote}
宋氏积弱，元兵南伐，势若摧枯：列郡土崩，不降则溃。其有孤城介立，血战经年，洎行在失守，三宫北迁，而焚诏斩使，勇气弥厉，忠胜于张巡，守坚于墨翟，则李庭芝乘城百战，国亡与亡也。当明季世，流寇滔天，南都草创，奸人在朝，方镇擅命，国势殆哉，不可为矣！其有上匡暗主，下抚骄将，内揽群策，外抗天兵，鞠躬尽力，死而后已，则史可法效命封疆，终为社稷臣也。
\end{BookQuote}

扬州民众在这两次战争中的命运如何，史有明载。先看李庭芝坚守下的扬州民众。史载，元将阿术率兵于南宋德佑元年（1275）初夏起进攻扬州，人间惨剧随之而至：8

\begin{BookQuote}
（元兵）久攻不拔，乃筑长围困之。冬，城中食尽，死者满道。明年二月，饥益甚，赴濠水死者日数百。道有死者，众争割，啖之立尽…。庭芝犹括民间粟
\end{BookQuote}

\vfill
\begin{BookNotes}
\BookNote{7.}{《宋史·李庭芝传》。}
\BookNote{8.}{同上。}
\end{BookNotes}

% 原书第 374 页（PDF 第 387 页）
\BookPage{374}{387}
以给兵。粟尽，令官人出粟。又尽，令将校出粟，杂牛皮、曲蘖以给之。兵有烹子而食者。犹日出苦战。

再看史可法抗清扬州城破后的民众。清顺治二年（1645）四月廿五日城破，清军屠城，史称「扬州十日」，死者逾八十万余:9

念五日，天渐暮，敌兵杀人声已彻门外。城中四周火起，近者十余处，远者不计其数。赤光相映如雷电，辟卜声轰耳不绝。又隐隐闻击楚声。哀顾断续，惨不可状。初二日，谕各寺院僧人焚化积尸。而寺院中藏匿妇女亦复不少，亦有惊饿死者。查焚尸簿载，其数约计八十万余。其落井投河，闭户自焚及深入自缢者不与焉。初三日，出示发赈。其往来负载者俱焦头烂额，断臂折胫，刀痕遍体，血渍成块，满面如烛泪成行。碎烂鹑衣，腥秽触鼻。人扶一杖，挟一蒲袋，正如神庙中窜狱冤鬼。初四日，天始霁，道路积尸既经积雨暴涨，而青皮如蒙鼓，血肉内溃，秽臭逼人。复经日炙，其气愈甚。前后左右，处处焚灼，腥闻百里。盖此百万生灵，一朝横死，虽天地鬼神，不能不为之愁惨也。

在邗城建成以来的两千多年历史中，频繁发生的内战与外（族入侵）战，多次给扬州城及其居民带来类似上述的毁灭性灾难，且看史载数例：

南朝刘宋大明三年(459)，皇室内订，明帝遣沈庆之攻广陵，自四月至七月，城破，明帝命屠城：10沈庆之请自五尺以下全之，其余男子皆死，女子以为军赏，犹杀三千余人。

南朝萧梁大宝元年（550），叛将侯景攻广陵：11

克之，执（守将）祖皓，缚而射之，箭遍体，然后车裂以徇。城中无少长皆埋

\vfill
\begin{BookNotes}
\BookNote{9.}{（清）王秀楚，《扬州十日记》，网上。}
\BookNote{10.}{《资治通鉴》。}
\BookNote{11.}{同上。}
\end{BookNotes}

% 原书第 375 页（PDF 第 388 页）
\BookPage{375}{388}
之于地，驰马射而杀之。

唐光启三年（887），地方军阀秦彦、毕师铎、孙儒、杨行密等混战，扬州民众受害之惨，不忍足闻：12杨行密围广陵且半年，…城中无食，斗米直钱伍十缗。草根木实皆尽，以堇泥为饼食之。饿死者太半。宣军（孙孺军）掠人，诣市卖之。驱缚屠割如羊豕，讫无一声。积尸流血，满于坊市…(城破后）遗民才数百家，饥羸无复人状。

后周显德二年（955），南唐中宗李璟避战，下令扬州徙民空城，居民蒙受巨大损失：13

知东都（扬州）必不守，遣使焚其官私庐舍，徙其民于江南。

南宋建炎三年（1129）正月，金兵南下攻宋，扬州又在劫难逃：14

二月，高宗渡江南逃，金兵入扬州，旋烧城北撤，居民十余万人，一部分死于江中，余悉被俘。元末，地方军阀混战，扬州残破，十室九空:15

（朱元璋军）按籍城中，仅余（土着）十八家。

南朝刘宋诗人鲍照（414-466）把战争残破后景象凄凉的广陵城称为「芜城」，并在其名作〈芜城赋〉中这样描述:16

木魅山鬼，野鼠城狐，风嚎雨啸，昏见晨趋。饥鹰厉吻，寒鸱吓雏…孤蓬自振，惊沙坐飞；灌木杳而无际，丛薄纷其相依；通池既已夷，峻隅又已颓。直视千里外，唯见起黄埃…边风急兮城上寒，井径灭兮邱陇残！

他是在前述沈庆之屠城不久之后到的扬州，汉晋以来繁华似锦的广陵城，已经荡然无存，眼前展现的是一座荒凉无比的荒芜之城一芜城。

\vfill
\begin{BookNotes}
\BookNote{12.}{《资治通鉴》。}
\BookNote{13.}{同上。}
\BookNote{14.}{沈起炜，《中国历史大事年表》，上海辞书，1983。}
\BookNote{15.}{《明太祖实录》。}
\BookNote{16.}{《昭明文选》，卷2。}
\end{BookNotes}

% 原书第 376 页（PDF 第 389 页）
\BookPage{376}{389}
乡贤郑板桥（1693-1766）曾云扬州「千年征伐百余次」17，虽属诗语，大致近之。故以上所引的战祸实例，实在是挂一漏万。汪文〈广陵对〉何以对此不及一语，似亦不难揣度。扬州民众历代身受兵祸之惨烈，宁有超过「扬州十日」者？身为邑人学者，于此宁有不闻不知之理？但身为清代臣民，面对官员公开谈史，自然要讳言本朝之恶。在文字狱恐怖之下，能将史可法上榜，已属非易。古今同理，不必苛求。

前人汪氏笔下之阙，后人郭沫若于史公祠题联补之：

\begin{BookQuote}
骑鹤楼头难忘十日，梅花岭畔共仰千秋。
\end{BookQuote}

上联悼念死难民众，下联礼敬殉国英雄。无一缺位，理当如此。

「难忘」与「共仰」，在五千年中华民族历史中，在九百多万平方公里中国大地上，扬州作为英雄城是受之无愧的。但是一定要记住，扬州人的遭遇不过是国人遭遇的一个缩影；而且无论苦难与忠义，毕竟都是既往陈事。我们衷心希望这种历史到此为止。今天之所以翻笈重温，意在以史为鉴而已：以人为本，珍惜和平；远离战争，无论内外；己所不欲，勿施于人。随着国力的增强，尤其要警惕动辄说「No」之风对国家和平发展及民众生命安全可能造成的严重危害。扬州作为英雄城和战争的巨大牺牲者，有权力、有责任为此大声疾呼。如有可能，在史公祠畔修建一座和平博物馆，以扬州遭受战祸之惨痛史实，教育后人珍惜和平、远离战争，不让历史悲剧重演，则善莫大焉。

当然，今后战争能否避免，历史悲剧会不会再度重演，并不取决于我们的善良愿望。在全球化的今天，只有站在全人类的高度，顺应时代潮流，从源头上切实解决好引发内战和外战的各种社会政治与经济问题，和平与和谐才不是一种奢望。在下面第（三）节中，我们还将就此话题稍作延伸讨论。

\vfill
\begin{BookNotes}
\BookNote{17.}{（清）郑板桥，〈扬州〉。}
\end{BookNotes}

% 原书第 377 页（PDF 第 390 页）
\BookPage{377}{390}
\BookSubheading{（二）繁华之都的背后}

古人云:「靠山吃山，靠水吃水。」今人说:「要致富，先修路。」古扬州的地理条件，堪称两者兼备，只要「天下无事」，理应富甲四方。诚如〈广陵对〉所言，扬州滨临东海，靠水吃水，「煮海为盐」，堪称物华天宝，得天独厚。至于「路」，南傍长江，北望淮河，也是得天独厚；另赖天赐「水利专家」隋炀帝，修成南北大运河，致使扬州坐收南北东西水运枢纽之利，成为交通与商贸中心，最终如〈广陵对〉所言，扬州「川渠所引，百货通焉，利尽四海」。

扬州的富庶繁荣程度，在唐则「扬一益二」，全国称冠；在明朝嘉万至明亡以及清朝康乾全胜时期，亦足比灵斯唐朝。请看清代扬州富商之多之富:18

\begin{BookQuote}
四方豪商大贾，鳞集麋至。侨寄户居者，不下数十万。…富者以千万计，百万以下者皆谓之小商。
\end{BookQuote}

盐业是清政府财政收入的主要来源之一，史载：19

\begin{BookQuote}
盐课居赋税之半，两淮盐课又居天下之半…损益盈虚，动关国计。
\end{BookQuote}

所以，麝集于扬州的豪商大贾，多为盐商。

盐业利大，商人乐营；盐课高收，政府乐管；兼以交通便捷，扬州焉有不富、不荣之理。古扬州富庶繁荣的标志是市场活跃、百业兴旺：建筑，园林，花卉，服饰，古玩，摆设，餐饮，书画，印刷，戏曲，教育，医药，美容，洗浴…等等，规模档次，不少都堪称当时国内一流，驰名四方。

市场活跃与百业兴旺，自然就业机会就多，谋生就比较容易，因而居民生活水准比之经济落后地区想必也会高出一筹。比如花卉，郑板桥说扬州「十里栽花算种田」20，可见从业者之多与消费市场之大。花卉业者的收入水准想必会高于种田不少。

\vfill
\begin{BookNotes}
\BookNote{18.}{《清朝野史大观》。}
\BookNote{19.}{（清)《两淮盐法志》。}
\BookNote{20.}{(清）郑板桥，《扬州》。}
\end{BookNotes}

% 原书第 378 页（PDF 第 391 页）
\BookPage{378}{391}
经济繁荣当然值得称道。但在古代社会缺乏收入调节机制的情况下，经济愈繁荣，贫富差距反而愈大，致使社会畸形变态加剧，社会和谐程度因之较经济落后地区可能不高反低。

扬州巨富当然首推盐商。其中固然有富而好学、施财济贫、热心公益者，如马曰琯、程梦星等；但总体而言，富埒王侯、显富摆阔、穷奢极欲、精神变态者不在少数，令人侧目：21

\begin{BookQuote}
扬州盐务，竞尚奢丽。婚嫁丧葬，堂室饮食，衣服舆马，动辄费数十万。有某姓者，每食，苞人备席十数类。临食时，夫妇并坐堂上，侍者抬席置于前，自茶面荤素等色，凡不食者摇其头，侍者审色，则更易其他类。或好马，蓄马数百，每马日费数十金。朝自内出城，暮自城外入。五花灿着，
\end{BookQuote}

观者目炫。

\begin{BookQuote}
或好兰，自门以至于内室，置兰殆遍。或以木作裸体妇人，动以机关，置诸斋阁，往往座客为之惊避。其后起之家，更有足异者：有欲以万金一时费去者，门下客以金尽买金箔，载至金山塔上，向风扬之，顷刻而散。沿江草树之间，不可收复。又有三千金尽买苏州不倒翁，流于水中，波为之塞。有喜美者，自司暗以至灶婢，皆选十数龄清秀之辈，或反之而极，尽用奇丑者，自镜子以为不称，毁其面以酱敷之，暴于日中。有好大者，以铜为溺器，高五六尺。夜欲溺，起就之。一时争奇斗异，不可胜记。
\end{BookQuote}

文中所谓「后起之家」，多指暴发户，钱来得太容易，精神愈加变态，行为愈显乖戾。其所作所为，照常理自然难以理解。其实，此辈代有传人，于今即遍地之X老板是也。以今视古，解亦不难。

古代经济繁荣之地，贫富悬殊导致社会失谐的丑恶现象之一是青楼倡盛。古扬州不但不例外，并且驰名全国，甚至在明清时期因之被称为「青楼都会」。22

\vfill
\begin{BookNotes}
\BookNote{21.}{（清）李斗，《扬州画舫录》。}
\BookNote{22.}{王振忠，〈明清徽商与扬州城市文化的特征和地位〉，载于《扬州研究》，台北版，1996。}
\end{BookNotes}

% 原书第 379 页（PDF 第 392 页）
\BookPage{379}{392}
\noindent 唐代青楼的盛况，由唐人诗文可窥一斑：扬州胜地也。每重城向夕，娼楼之上，常有绛纱灯万数，辉罗耀列空中。九里三十步街中，珠翠塡咽，邈若仙境。23十里长街市井连，月明桥上看神仙。24当年人未识干戈，处处青楼夜夜歌。25夜市千灯照碧云，高楼红袖客纷纷。如今不是时平日，犹自笙歌彻晓闻。26

清朝乾隆时期同样青楼倡盛。乡贤李斗所述的扬州青楼盛况，与唐代何其相似乃尔：27

每夕燃灯数万，粉黛绮罗甲天下。

其时青楼爆门者，多为富商，士大夫亦不落后:28

士大夫俱尚豪华，而尤喜狎邪之游。在江宁则秦淮河上，在苏州则虎丘山塘，在扬州则天宁门外之平山堂。画船箫鼓，殆无虚日。

唐朝不少著名文士都在扬州浪迹过烟花柳巷，并赋诗吟咏，其中不乏名篇，例如杜牧的〈遣怀〉，就直捷了当道出自己是青楼常客，丝毫不加隐讳：

落魄江湖载酒行，楚腰纤细掌中轻。十年一觉扬州梦，赢得青楼薄幸名。

《唐诗别传》对此加以证实：牧在扬州每夕为狎邪游，所至成欢，无不会意，如是者数年。

但，杜牧〈赠别〉诗中的妓女不仅被写成美貌绝顶，而且在与狎客离别时的表现，如同挚爱的情侣般难解难分：娉娉袅袅十三余，豆蔻梢头二月初。春风十里扬州路，卷上珠帘总不如。

\vfill
\begin{BookNotes}
\BookNote{23.}{（唐）杜牧，《樊川诗集》。}
\BookNote{24.}{（唐）张祜，〈纵游淮南〉。}
\BookNote{25.}{（唐）韦庄，〈过扬州〉。}
\BookNote{26.}{（唐）王建，〈夜看扬州市〉。}
\BookNote{27.}{（清）李斗，《扬州画舫录》。}
\BookNote{28.}{（清）钱泳，《履园丛话》，卷7。}
\end{BookNotes}

% 原书第 380 页（PDF 第 393 页）
\BookPage{380}{393}
\begin{BookQuote}
多情总觉反无情，唯觉尊前笑不成。蜡烛有心还惜别，替人垂泪到天明。
\end{BookQuote}

在诗人徐凝的名作《忆扬州》中，类似的描写尤为动人：

\begin{BookQuote}
萧娘脸上难胜泪，桃叶眉头易得愁。天下三分明月夜，二分无赖是扬州。
\end{BookQuote}

诗中娉娉嫘嫘者、惜别垂泪不绝者、萧娘和桃叶都是妓女。妓女与嫖客之间，在嫖客杜牧与徐凝的笔下，竟是如此温情脉脉、情意缠绵，符合实际吗？或许有之，但绝非多数。妓女的真实感受，请听《老残游记》里妓女翠环是怎么说的：

\begin{BookQuote}
至于说姐儿怎么跟他好，恩情怎么重，我有一回发了傻性子，去问了问，那个姐儿说：「他住了一夜，就麻烦了一夜。天明问他要讨个两数银子的体己，他就抹下脸来直着脖儿梗，乱嚷道：「我正账昨儿晚上就开发了，还要什么体己钱？」那姐儿呢，再三央告说：「正账的钱呢，店里伙计扣一分，掌柜的又扣一分，剩下的全是领家的妈拿去，一个钱也放不出来。俺们的胭脂花粉，跟身上穿的小衣裳，都是自己钱买。」再三央告着，他给了二百钱一个小串子，望地下一摔，还要撅着嘴说：「你们这些强盗婊子，真不是东西！混账王八蛋！」你想有恩情没有？因此，我想，做诗这件事是很没有意思的，不过造些
\end{BookQuote}

谣言罢了。

她说的「做诗」，是指嫖客眠花宿柳之作，也就是与上述杜牧、徐凝相似的那类作品。嫖客于狎邪游之后赋诗，或自娱或炫耀他人，在当时社会常常被视为风流韵事，无须忌讳，一般与什么「恩情」根本无涉。所以妓女翠环鄙夷这类诗作不过是「造些谣言」罢了。如此评论，精辟透顶，振聋发聩。常见的嫖客的真实「恩情」在那里？我们从翠环话里，听的还不够清楚吗？

古代经济繁荣之地，贫富悬殊导致社会失谐的另一丑恶现象是在狎妓之外，富商和士大夫们还要纳妾。《儒林外史》中的扬州盐商宋为富说：

\begin{BookQuote}
我们总商人家，一年至少也娶七八个妾。
\end{BookQuote}

纳妾要花钱，宋老板有的是银子。他想骗娶一位贡生女儿做妾，出银五百两。

郑板桥的漂亮如夫人贫家女子饶五，出嫁前差点被一富商买走，出价是七百


% 原书第 381 页（PDF 第 394 页）
\BookPage{381}{394}
\noindent 金。出道前在扬州打工卖字卖画的一介寒士郑板桥，自云当时生活状况是「落拓扬州一蔽裘，绿杨萧寺几淹留」、「乞食山僧庙，缝衣歌妓家」，怎么娶得起妾？据说是靠富商程梦星的资助。板桥终生仅此一妾，夫妻感情不错，也算是饶女士造化。29

明清时期，一度流行所谓「要娶小，扬州讨」之说，扬州俨然成为「青楼都会」。[参见注10]

尽人皆知，妓女与媵妾的命运无比悲惨，但凡日子能过谁乐为之！郑板桥云扬州「千家养女先教曲」30，是像我们当今城市家庭儿童那样为升学加分而上音乐班吗？当然不是！说来心酸，那是蓬户人家为女儿将来或娼或妾卖个好价，从小就教孩子学唱。青楼女子、富家媵妾的来源，主要就在于此。

古扬州繁华时期富商、士大夫多，狎妓、纳妾市场需求大。穷人更多，「货源」更充足。当然穷人中宁可饿死，亦不操丑业，占多数。而为各种不得已因素所迫，总有一些人沦为娼妓、媵妾，在当时，这种丑恶现象谁也无法改变。

古扬州竟因此而蒙羞。谁之过？繁荣之过？繁荣本身，大好事也，何过之有！实实在在是当时不良社会制度之过！

难忘西元1949年！新政权建立后的最大善政之一，是立即在全国取缔妓院、解放妓女、严禁纳妾，从此卖淫嫖娼和纳妾现象绝迹，社会面貌为之一新。大快人心事也！

然而，不幸的是，随着国家经济迅速发展，种种旧时代丑恶现象包括卖淫嫖娼，又在全国各地不同程度地死灰复燃。历史与现实，竟然又在这里交汇，岂偶然哉？与战争起因一样，其根源同样也在社会政治与经济制度。

\vfill
\begin{BookNotes}
\BookNote{29.}{《扬州八怪传》，网上。}
\BookNote{30.}{（清）郑板桥，〈扬州〉。}
\end{BookNotes}

% 原书第 382 页（PDF 第 395 页）
\BookPage{382}{395}
\BookSubheading{（三）年年红药为谁开}

行文至此，作为和平与美好的象征，我们不禁想到扬州的名花红芍药。

「维士与女，伊其相谑，赠之以芍药」31。在两三千年前的郑国，男女相亲相爱，即以芍药相赠，足见此花形美而色艳，传情致意，招人喜爱。后世广泛种植，至北宋时期，扬州芍药名满天下，与洛阳牡丹并列百花之冠。北宋文学家欧阳修守扬时，曾题诗加以赞美：

\begin{BookQuote}
琼花芍药世无伦，偶不题诗便怨人。曾向无双亭下醉，自知不负广陵春。
\end{BookQuote}

人赏花，花解诗。人花相守，多么和谐，多么美好。可惜的是，当时人世间有闲暇逸致持酒品花如欧阳公者，能有几人？欧阳公辞世不过五十多年，金兵南侵，北宋灭亡，天下大乱，人命难保，遑论赏花。

大约距欧阳公赋诗一百年后，二十一岁的南宋词人姜夔首次路过扬州，特地慕名前往游览胜地二十四桥游玩。对这里，他从小就因杜牧的动人诗句「二十四桥明月夜，玉人何处教吹箫」而一直向往不已。可置身其地，眼前的荒芜凄凉景象令他十分震惊与惆怅，情不自禁地喟然咏叹道：32

\begin{BookQuote}
二十四桥仍在，波心荡，冷月无声。
\end{BookQuote}

不见玉人，不闻箫声；寒水默默自荡，月色幽茫凄迷；杂草丛生，了无花影；唯一堪以自慰的是「桥仍在」。

他来的不是时候。这时的扬州，既非杜牧青楼得意时的唐代首富大都会，也不是欧阳修「坐花戴月」时的北宋腹地繁华之邦，而是两度惨遭金兵残破后的南宋边陲戍城。这座城市蒙受的最近一次战祸劫难，就发生在十五年前金主完颜亮的南侵。展现在他眼前的扬州，可能与七百年前鲍照眼中的「芜城」相差无几，几乎就是一座空城：33

\vfill
\begin{BookNotes}
\BookNote{31.}{《诗经·郑凤·溱淆》。}
\BookNote{32.}{（南宋）姜夔，〈扬州慢〉。}
\BookNote{33.}{（南宋）姜夔，〈扬州慢〉。}
\end{BookNotes}

% 原书第 383 页（PDF 第 396 页）
\BookPage{383}{396}
\begin{BookQuote}
淮左名都，竹西佳处，解鞍少驻初程。过春风十里，尽荠麦青青。自胡马窥江去后，废池乔木，犹厌言兵。渐黄昏，清角吹寒，都在空城。
\end{BookQuote}

比姜夔晚四百多年的明清诗人吴伟业，在「扬州十日」十年之后来到满目疮痍的扬州城，感念故国故地，心情沉痛地咏叹道：34

\begin{BookQuote}
十载西风空白骨，廿桥明月自朱楼。
\end{BookQuote}

「空」者，徒劳无功也：近百万人的牺牲，累累西风白骨，没能改变亡国的命运；「自」者，物是人非也:二十四桥仍在，朱楼明月谁赏？

鲍照、姜夔和吴伟业等人的诗文〈芜城赋〉、〈扬州慢〉、〈扬州〉，不仅给战争残破后的扬州留下凄凉憔悴的身影，更给后人留下不少值得永远思考的问题。

鲍照面对芜城扬州，抚今思昔，感慨兴亡荣枯之无常，发出惊天一问：35天道如何？

「天道如何？」问的是国家为何有兴有亡、有治有乱。这个问题，出题者鲍照本人并未给出明确的回答，只是凭感觉说:「吞恨者多」36。

亲罹外族入侵战乱和深怀亡国之痛的吴伟业，相信「天道如何」之问也会久久回荡在他的脑海。他在〈扬州〉诗里对此似乎未着一笔。但，一个「空」，一个「自」，给后世留下多大耐人寻味的思索空间！

不曾料到，熟谂历代史事的汪中，在〈广陵对〉中对于类似鲍照「天道如何」的问题竟然作出如下的回答：

\begin{BookQuote}
至若政事法理、经纬乎民生…一至之行，盖以千百计，非国家之所以废兴存亡
\end{BookQuote}

者，则皆略之。

国家废兴存亡，怎么会与「政事法理、经纬乎民生」无关？国家治理，悠悠万事，宁有重于「民生经纬」与「政事法理」者？「百代兴亡朝复暮，东风吹倒前朝

\vfill
\begin{BookNotes}
\BookNote{34.}{（清）吴伟业，《扬州》。}
\BookNote{35.}{（南朝宋）鲍照，（芜城赋〉。}
\BookNote{36.}{同上。}
\end{BookNotes}

% 原书第 384 页（PDF 第 397 页）
\BookPage{384}{397}
\noindent 树」37。自秦以来两千多年，内战外战，天下大乱、王朝更迭，兴勃亡忽，举其大者不下二十余番，几乎无一不因贫富悬殊、民不聊生、社会不公、政治腐败所致。汪氏可谓误矣。

在繁华不再、满目萧索的空城扬州，在人迹罕至的二十四桥一红药桥旁边，多情诗人姜夔惆怅之余，不禁为今后年复一年的桥边红药悬念不已:38

\begin{BookQuote}
念桥边红药，年年知为谁生？
\end{BookQuote}

在南宋偏安江淮之南一隅、国势日蹙的情况下，诗人简直不敢想像扬州会很快重现往日辉煌。所以他的问题像是凌空一箭射向未来：桥边美丽的红芍药啊，以后你年年为谁而开？年年花开又有谁来欣赏？

谁？谁？太尖锐了！姜夔心中的「谁」，可能指的就是他的同类文士，比如唐朝的杜牧，这由他的词文可鉴：39

\begin{BookQuote}
杜郎俊赏，算而今，重到须惊。纵豆蔻词工，青楼梦好，难赋深情。
\end{BookQuote}

尽管鲍照和姜夔的「天道如何」、「桥边红药，年年知为谁生」之问，是问向战争残破后的「芜城」「空城」扬州，但只要人类尚处追求「世界大同」的历史进程之中，我们就不得不时时要去考虑对之如何加以回答，即使是在经济繁荣的和平时期。

比如，在唐、清经济繁荣的和平时期，二十四桥美丽的红芍药，年年仅仅为杜牧和姜夔等士大夫以及富人宋为富等少数人开放，他们是心满意足了，社会就会因此而和谐美满了吗？国家就会因此而永葆太平了吗？且看在扬州「乞食山僧庙」、「缝衣歌妓家」时的打工仔郑板桥，他眼里见到的可不是含情脉脉的红芍药，而是「廿四桥边草径荒」40，生计艰难使他对社会怨气满腹:「尽把黄金通显要，惟余白眼到清贫」。41更不要说在家待价卖身的贫女饶五们了。杜牧、姜夔、宋为富、

\vfill
\begin{BookNotes}
\BookNote{37.}{（清）吴敬梓，《儒林外史》。}
\BookNote{38.}{（南宋）姜夔，扬州慢）。}
\BookNote{39.}{同上。}
\BookNote{40.}{（清）郑板桥，〈扬州〉。}
\BookNote{41.}{同上。}
\end{BookNotes}

% 原书第 385 页（PDF 第 398 页）
\BookPage{385}{398}
\noindent 郑板桥、饶五…，各有各的天地，各有各的诉求，只有国家的政经制度能够确保将之一一协调到位，社会方能和谐美满，国家方能永葆太平，美丽的红芍药才会为社会每一对情侣鲜艳绽放。

遗憾的是，自秦至清，「百代都行秦政治」，而「秦政治」恰恰与社会协调的政经要求背道而驰，所以无一能逃兴勃亡忽的「周期率」:「百代兴亡朝复暮」。

今人问政，史鉴俱在，可不慎欤？

谨以此文祝福乡邦扬州:年年红药为众开，居者怡而远者慕，和谐美好可人，「人生只爱扬州住」，「人生只合扬州死」，岂不休哉！


% 原书第 386 页（PDF 第 399 页）
\BookPage{386}{399}
